\section{C. Extension and descent}
\subsection{Nonreduced bases: two opposite examples}

\begin{proposition}[C-RD3: nilpotent invariance when $p$ is locally nilpotent]\label{prop:C-RD3}
\status{Proved.}
Let $S$ be locally noetherian, suppose $p$ is locally nilpotent on $S$, and
let $U\subset S$ be open. Write $S_0=S_{\mathrm{red}}$ and
$U_0=U\times_SS_0$. A triple $(A_U,M,\varphi)$ lies in the essential image
of $F_S$ if and only if its reduction lies in the essential image of
$F_{S_0}$. If $F_{S_0}$ is fully faithful, so is $F_S$. In particular,
if $F_{S_0}$ is an equivalence, then $F_S$ is an equivalence.
\end{proposition}
\begin{proof}
The forward implication is restriction. Conversely let $A_0/S_0$
realize the reduced triple, with the given isomorphism
$A_0[p^\infty]\simeq M|_{S_0}$. Locally on the noetherian scheme $S$,
the nilradical is nilpotent. The Serre--Tate equivalence
\cite[Theorem~1.2.1]{Katz1981SerreTate} constructs an abelian scheme
$A/S$ lifting $A_0$ and realizing $M$. Its restriction to $U$ and $A_U$
have identical reduction and identical $p$-divisible group, with the
specified identification. The full faithfulness in the same theorem
provides their unique compatible isomorphism. These local constructions
glue by that uniqueness. For morphisms, first descend the reduced pair
using full faithfulness of $F_{S_0}$ and then apply full faithfulness of
Serre--Tate. This proves every assertion.
\end{proof}

\begin{corollary}[C-RD4: the proposed reducedness criterion is false]\label{cor:C-RD4}
\status{Refuted, with a positive counterexample.}
For $k=\overline{\mathbf F}_p$, let
\[
 S=\operatorname{Spec}k[[t]][\epsilon]/(\epsilon^2),\qquad U=D(t).
\]
The hypotheses a) and b) of 100:3 hold, the local ring at the point of
$T=S\setminus U$ is nonreduced, and nevertheless $F_S$ is an equivalence
in all relative dimensions.
\end{corollary}
\begin{proof}
The finite surjective map $S_0=\operatorname{Spec}k[[t]]\to S$ supplies
a), since $S_0$ is normal and geometrically unibranch. The etale topoi
are unaffected by the nilpotent thickening. A strictly henselian trait
has connected punctured spectrum, so its local etale cohomology with
locally constant coefficients vanishes in degrees zero and one. This
is $\operatorname{prof}^{\mathrm{et}}_T(S)\geq2$; it verifies b).

Over the trait $S_0$, the existence of an extension of $A_{U_0}[p^\infty]$
implies good reduction of $A_{U_0}$ by de Jong's good-reduction criterion
\cite[Theorem~3, citing Theorem~2.5 of the original paper]{deJong1998ICM}.
The given identification with $M_0$ extends uniquely, by the extension
theorem for homomorphisms of Barsotti--Tate groups
\cite[Theorem~2]{deJong1998ICM}. Morphisms of the abelian schemes over a
trait extend by the Neron mapping property. Thus $F_{S_0}$ is an
equivalence. Proposition~\ref{prop:C-RD3} gives the assertion for $S$.
The nonzero element $\epsilon$ proves nonreducedness at $T$.
\end{proof}

\begin{proposition}[C-RD5: an explicit failure of 100:3(2)]\label{prop:C-RD5}
\status{Proved counterexample; R4 rechecked.}
Let
\[
 S=\operatorname{Spec}\mathbf Z_5[\epsilon]/(\epsilon^2),\quad U=D(5),
 \quad E/\mathbf Z_5:\ y^2=x^3+x+1.
\]
Put $M=E[5^\infty]\times_{\mathbf Z_5}S$, and over $U$ take
\[
 A_U:\quad y^2=x^3+x+(1+\epsilon).
\]
There is a unique isomorphism $M|_U\simeq A_U[5^\infty]$ reducing to the
identity modulo $\epsilon$. This triple does not extend to an abelian
scheme over $S$, although a) and b) of 100:3 hold.
\end{proposition}
\begin{proof}
Here $\Delta(E)=-496$ is a $5$-adic unit, so $E$ is an elliptic scheme.
The two generic $5$-divisible groups are etale; invariance of finite
etale schemes under nilpotent thickenings gives the stated compatible
isomorphisms at all levels. Hypothesis a) holds via
$\operatorname{Spec}\mathbf Z_5\to S$, and b) follows from the connected
punctured strict localization of this trait, as above.

If $A/S$ realized the triple, its reduction modulo $\epsilon$ would be
$E$: both are good models of the same elliptic curve over $\mathbf Q_5$.
For every $n\geq1$, apply Serre--Tate over
$(\mathbf Z_5/5^n)[\epsilon]/(\epsilon^2)$. The specified constant
$5$-divisible-group deformation forces $A$ modulo $5^n$ to be the
constant deformation of $E$ modulo $5^n$. Hence
$j(A)-j(E)$ lies in $\bigcap_n5^n\mathbf Z_5[\epsilon]=0$.
On the other hand, the generic fibre prescribed above has
\[
 j(A_U)=\frac{6912}{31+54\epsilon}
 =\frac{6912}{31}-\frac{373248}{961}\epsilon,
\]
whose $\epsilon$ coefficient is nonzero. This contradiction proves the
claim. It also identifies exactly the failure of the proposed gluing
at 100:7: equality of the generic etale $p$-divisible groups supplies no
isomorphism between these two characteristic-zero infinitesimal
abelian-scheme deformations.
\end{proof}

\subsection{The normal extension theorem, including nonexcellent bases}

\begin{theorem}[C-RD6: the normal case]\label{thm:C-RD6}
\status{Proved by the checked normal-base extension theorem.}
Let $S$ be a locally noetherian normal scheme, let $U\subset S$ be a dense
open, and suppose every point of $S\setminus U$ has residue
characteristic $p$. The functor $F$ of Part C is an equivalence.
Neither excellence nor finiteness of normalizations of modifications of
$S$ is needed. In particular the normal case of C.2 holds without
hypotheses a) or b), apart from the indicated density.
\end{theorem}
\begin{proof}
Work on one connected component, which is integral. For a triple
$(A_U,M,\varphi)$, \cite[Corollary~6.12]{KrishnamoorthyPal2021RankTwo}
produces an abelian scheme $A/S$ extending $A_U$. This theorem has
exactly the hypotheses displayed here: locally noetherian and normal,
with the missing locus in characteristic $p$. It is not a purity
statement requiring a regular or unramified base.

A normal noetherian integral scheme is a Krull scheme. Apply
\cite[Theorem~6.3]{KrishnamoorthyPal2021RankTwo} to the generic-fibre
isomorphism $A[p^\infty]\simeq M$ induced by $\varphi$. That theorem
states the full faithfulness of restriction of Barsotti--Tate groups
from an integral Krull scheme to its generic point. Apply it to the
isomorphism and its inverse; their composites are identities by
faithfulness. This gives the required isomorphism over $S$. Its
restriction to $U$ is the original one, again by faithfulness. Thus the
entire triple, including its specified identification, is realized.
For full faithfulness, a morphism of triples gives a homomorphism of
BT groups over $S$ which is algebraic over the dense open $U$. Apply
\cite[Proposition~1.2, p.64]{Grothendieck1966Hom}: for every associated
point of the reduced scheme $S$, the required algebraicity hypothesis
holds at that generic point itself. The proposition produces the unique
abelian-scheme homomorphism inducing the prescribed BT map. Its
restriction to $U$ is the original homomorphism by faithfulness.
All constructions agree on overlaps by uniqueness and hence glue.

For clarity, the nonexcellent point is explicitly addressed in the
source's proof of Corollary~6.12. It takes the closure of the graph of a
map to a fine polarized moduli scheme, then passes to its normalization,
which need not be finite or noetherian. The normalization is an integral
Krull scheme, so Theorem~6.3 applies. The proof uses the Mori--Nagata
statements cited there as EGA~$0_{\mathrm{IV}}$, 23.2.6--23.2.8 to
control the reduced fibres, then Corollary~6.10 and Proposition~6.11 to
show that the graph has zero-dimensional fibres. This is precisely the
normalization issue that cannot be omitted by imposing an unstated
excellence assumption.
\end{proof}

\begin{remark}[An additional gap in the inherited spreading argument]\label{auto:C:4b7cbcb1fe87}
\status{Proved under stated hypotheses}

The proof in the supplied reading of 100:15--17 invokes EGA~IV, 8.8.2
to spread both an isomorphism of abelian schemes and an isomorphism of
$p$-divisible groups from a local scheme to a neighbourhood. The first
invocation is justified by finite presentation. The second is an
infinite compatible system and does not follow from that theorem.
One needs a separate finite-level determination or an algebraicity-locus
argument. Consequently that inherited argument is not used to remove
normality in Theorem~\ref{thm:C-RD6}.
\end{remark}

\subsection{A reduced counterexample with nonfinite normalization}

\begin{lemma}[A differential subring with an analytically ramified completion]
\label{lem:C-RD11-ring}
\status{Proved under stated hypotheses}
Let $V$ be a complete discrete valuation ring of mixed characteristic
$(0,p)$ with uniformizer $p$, fraction field $K$ and residue field $k$.
Let $C\subset K$ be a subfield such that $C\cap V\to k$ is onto.
Let $J\in V$ be transcendental over $C$, let $F\subset K$ be an
algebraic extension of $C(J)$, and let $D:F\to F$ extend the derivation
which is zero on $C$ and sends $J$ to $1$. Put
\[
 V_0=F\cap V,\qquad R=\{a\in V_0:D(a)\in V_0\}.
\]
Then $R$ is a one-dimensional noetherian local domain, with fraction
field $F$ and residue field $k$. Its normalization is the DVR $V_0$,
which has uniformizer $p$ and residue field $k$; both the fraction-field
and residue-field extensions induced by $R\subset V_0$ are identities.
In particular $R$ is geometrically unibranch. Its normalization is not
finite, and the homomorphism
\[
 R\longrightarrow V[\epsilon]/(\epsilon^2),
 \qquad a\longmapsto a+\epsilon D(a)
                                                        \tag{C-RD11.1}
\]
identifies its maximal-ideal completion with
$V[\epsilon]/(\epsilon^2)$.
\end{lemma}
\begin{proof}
The ring $V_0$ is a valuation ring with value group $\mathbf Z$ and
uniformizer $p$, hence is a DVR. Every $f\in F$ can be multiplied by
a sufficiently large power of $p$ to make both $f$ and $D(f)$ integral;
therefore $F=R[1/p]=\operatorname{Frac}R$. Elements of $R$ of valuation
zero are units of $R$, since $D(a^{-1})=-a^{-2}D(a)$ is integral.
Thus $R$ is local, its maximal ideal is
$\mathfrak m=\{a\in R:v(a)>0\}$, and its residue field is $k$.

Choose $c\in C\cap V$ with the same residue as $J$ and put $z=J-c$.
Then $z\in\mathfrak m$ and $D(z)=1$. For $a\in\mathfrak m$, choose
$b\in C\cap V$ lifting the residue of $D(a)$. Both $a-bz$ and its
derivative are divisible by $p$, so $(a-bz)/p\in R$. Consequently
\[
 \mathfrak m=(p,z),\qquad \mathfrak m^2=p\mathfrak m,
 \qquad \mathfrak m^n=p^{n-1}\mathfrak m\quad(n\geq1). \tag{C-RD11.2}
\]
For the middle equality, observe that $z^2/p\in\mathfrak m$: its
valuation is positive and its derivative $2z/p$ is integral.

If $0\ne a\in R$ has valuation $r\geq1$ and
$b\in\mathfrak m^{2r}$, then
$v(b)\geq2r$, $v(D(b))\geq2r-1$ and $v(D(a))\geq0$. Hence
$b/a$ and
\[
 D(b/a)=\frac{aD(b)-bD(a)}{a^2}
\]
are integral. Thus $\mathfrak m^{2r}\subset aR$.
Since $\mathfrak m$ is finitely generated and $R/\mathfrak m=k$,
$R/aR$ is an artinian noetherian ring. Every nonzero ideal $I$ contains
such an $a$ and is finitely generated by lifting generators of $I/aR$.
This proves that $R$ is noetherian. Every nonzero prime contains a
power of $\mathfrak m$, so $\dim R=1$.

The residue-surjectivity assumption makes $C\cap V$ dense in $V$:
successively choose representatives of the residues of the remainders
after division by $p$. For $a\in V_0$, choose $N$ sufficiently large
that $p^Na^{p^N-1}D(a)$ is integral. Then $a^{p^N}\in R$, so $a$
is integral over $R$.
As $V_0$ is integrally closed, it is exactly the normalization. The
normalization has one point above the closed point and the same
residue field, which is the geometric-unibranch criterion.

Write $T=V[\epsilon]/(\epsilon^2)$ and
$\mathfrak n=(p,\epsilon)$. Formula (C-RD11.2) shows that the
inverse image of $\mathfrak n^n=p^nV+\epsilon p^{n-1}V$ under
(C-RD11.1) is exactly $\mathfrak m^n$. The image is dense: given
$a,b\in V$, approximate them by $a_0,b_0\in C\cap V$, and choose
$c_N\in C\cap V$ with $v(J-c_N)\geq N$. Then
\[
 r=a_0+b_0(J-c_N)\in R,\qquad D(r)=b_0,
\]
approximates $a+\epsilon b$ to the required order. Thus
$R/\mathfrak m^n\simeq T/\mathfrak n^n$ for every $n$, proving the
completion assertion.

Finally no nonzero $d\in R$ can satisfy $dV_0\subset R$. With $c_N$
as above put $u_N=(J-c_N)/p^N\in V_0$. Then
\[
 D(du_N)=D(d)u_N+d/p^N.
\]
For $N>v(d)$ the first term is integral and the second has negative
valuation, so $du_N\notin R$. A finite birational normalization
would have a nonzero conductor, by taking a common denominator for
finitely many generators. Therefore $V_0/R$ is not finite.
\end{proof}

\begin{theorem}[C-RD11: the reduced assertion fails without finite normalization]
\label{thm:C-RD11}
\status{Proved; independent audit of the construction recorded.}
There is a reduced, one-dimensional, geometrically unibranch,
noetherian local domain $R$ of mixed characteristic $(0,5)$, with
nonfinite normalization, for which hypotheses a) and b) of 100:3
hold, but the functor $F$ is not essentially surjective.
The failure already occurs for a split ordinary height-two BT group
and an elliptic generic fibre which itself has a good model over $R$.
Thus the assertion in C.2 for every reduced base is false.
\end{theorem}
\begin{proof}
Let $k=\overline{\mathbf F_5(t)}$ and $V=W(k)$.
Take the elliptic curve over $k$ with $j=t$. It is ordinary: in the
short model
\[
 Y^2=X^3-\frac{3t}{t-1728}X-\frac{2t}{t-1728},
\]
the coefficient of $X^4$ in the square of the cubic, namely twice
its $X$ coefficient, is nonzero in characteristic $5$.
Let $E^{\mathrm{can}}/V$ be its canonical ordinary lift and write
$J=j(E^{\mathrm{can}})$. The canonical split BT description is
\cite[Theorem~2.1]{Katz1981SerreTate}; in relative dimension one the
formal lift algebraizes using the relatively ample divisor given by
three times the zero section. We have
\[
 E^{\mathrm{can}}[5^\infty]\simeq
 \mu_{5^\infty}\oplus\mathbf Q_5/\mathbf Z_5.
\]
The residue $\bar J=t$ is transcendental over $\mathbf F_5$, so
$J$ is transcendental over $\mathbf Q$.

We describe the auxiliary subfield without a hidden transcendence
assumption. The field $k$ is countable, whereas every residue class
in $V$ is uncountable. Enumerate $k$ and choose a representative
$c_n\in V$ of its $n$th element which is transcendental over
$\mathbf Q(J,c_1,\ldots,c_{n-1})$. This is possible because the
algebraic closure of that countable field in $\operatorname{Frac}V$
is countable. Let $C=\mathbf Q(c_1,c_2,\ldots)$. Then $C\cap V\to k$
is onto, $C\cap V$ is dense in $V$, and $J$ is transcendental over $C$.

Since $J$ and $J-1728$ are units, the elliptic scheme
\[
 E_J:\quad Y^2=X^3-\frac{3J}{J-1728}X
                         -\frac{2J}{J-1728}             \tag{C-RD11.3}
\]
has $j$-invariant $J$ and discriminant
$1728^2J^2/(J-1728)^3\in V^\times$.
It is isomorphic over $V$ to $E^{\mathrm{can}}$: elliptic schemes
with the same $j$ away from $0,1728$ differ by a quadratic twist, and
every unit of the complete DVR $V$ is a square because $k$ is
algebraically closed and $2$ is invertible. Fix a compatible BT
isomorphism from $E_J[5^\infty]$ to the split group over $V$.

Each of its generic finite-level isomorphisms is defined over a
finite algebraic extension of $C(J)$ inside $\operatorname{Frac}V$:
the isomorphism scheme of the two finite etale group schemes over
$C(J)$ is finite etale. Let $F$ be the compositum of those countably
many fields. It is algebraic over $C(J)$ and the entire compatible
generic BT isomorphism is defined over $F$. Apply
Lemma~\ref{lem:C-RD11-ring} to this $F$ and the derivation $D(J)=1$,
$D(C)=0$, and obtain $R\subset F$.

Set $S=\operatorname{Spec}R$, $U=D(5)=\operatorname{Spec}F$, and
\[
 M=\mu_{5^\infty,R}\oplus(\mathbf Q_5/\mathbf Z_5)_R,
 \qquad A_U=(\text{equation (C-RD11.3)})_F.
\]
The coefficients of that equation belong to $F$; the chosen finite-level maps give the
required isomorphism on $U$. In fact $J\in R$ and $J-1728$ is a unit
of $R$, so the same equation is already an elliptic scheme over $R$.
Hypothesis a) holds with the identity cover, since $R$ is geometrically
unibranch. Hypothesis b) holds because the punctured strict local
spectrum is connected: a reduced geometrically unibranch local
domain has a reduced irreducible strict henselization, and in this
one-dimensional situation its punctured spectrum has a single
generic point. The closed point has characteristic $5$.

There is already an obstruction modulo $5$; completion is unnecessary
for the contradiction. The homomorphism (C-RD11.1), reduced modulo
$5$, induces
\[
 R/5R\simeq k[\epsilon]/(\epsilon^2),\qquad
 a\longmapsto\bar a+\epsilon\overline{D(a)}.            \tag{C-RD11.4}
\]
The kernel is $5R$ because $a,D(a)\in5V_0$ is equivalent to $a/5\in R$.
For surjectivity use representatives from $C\cap V$ and $z=J-c$,
which maps to $\epsilon$.

Suppose an abelian scheme $A/S$ realized this triple. It has relative
dimension one. Its generic $j$-invariant is $J$, and $R\hookrightarrow F$
is injective, so $j(A)=J\in R$. Under (C-RD11.4) this becomes
$t+\epsilon$. On the other hand, the reduction of $M$ to
$k[\epsilon]/(\epsilon^2)$ is the constant split BT group. If $A_0/k$
is the special fibre of $A$, Serre--Tate applied to the square-zero
thickening $k\hookrightarrow k[\epsilon]/(\epsilon^2)$ identifies
$A$ modulo $5$ with the constant deformation of $A_0$, because their
BT deformations are identified with the same constant group
\cite[Theorem~1.2.1]{Katz1981SerreTate}. Its $j$-invariant must
therefore be $t$. This contradicts $j(A\bmod5)=t+\epsilon$.
\end{proof}

\begin{remark}[Finite symbolic verification]\label{auto:C:3d23f9b8075d}
\status{Proved under stated hypotheses}

The exact symbolic certificate \path{verify_nonfinite_reduction.py}
checks the displayed discriminant, $j$-invariant, ordinary Hasse
coefficient and the nonzero dual-number coefficient. It makes no
claim to computationally instantiate the infinite field construction.
\end{remark}

\begin{remark}[Where the infinite system fails to descend]\label{auto:C:2b1cfc4567d4}
\status{Proved under stated hypotheses}

Over the normalization $V_0$, the chosen generic BT isomorphism
extends, since its finite-level coefficients were chosen inside
$F\cap V$. It does not descend to $R$. The normalization is integral
and radicial but not finite or proper of finite type, so this does
not contradict Proposition~\ref{prop:C-RD1}. It also supplies a
concrete reason that a finite-presentation limit theorem cannot be
applied to an entire BT isomorphism at once.
\end{remark}

\subsection{Morphisms of Barsotti--Tate groups and proper epimorphisms}

\begin{proposition}[C-RD1: proper epimorphisms]\label{prop:C-RD1}
\status{Proved.}
Let $S$ be a noetherian scheme, let $f:S'\to S$ be proper, and suppose
$\mathcal O_S\to f_*\mathcal O_{S'}$ is injective. Put $S''=S'\times_S S'$.
For any two $p$-divisible groups $G,H$ over $S$, the sequence
\[
 \operatorname{Hom}_S(G,H)\longrightarrow
 \operatorname{Hom}_{S'}(G_{S'},H_{S'})
 \rightrightarrows
 \operatorname{Hom}_{S''}(G_{S''},H_{S''})
\]
is an equalizer. Thus the corollary proposed at 100:25 is true. No hypothesis
that $S$ be reduced, normal, or excellent is needed.
\end{proposition}

\begin{proof}
We first prove the precise rigidity statement needed below. If $R$ is
noetherian, $J\subset R$ is nilpotent, and $G,H$ are $p$-divisible groups over
$R$, the reduction map
\[
 \operatorname{Hom}_R(G,H)\longrightarrow
 \operatorname{Hom}_{R/J}(G_{R/J},H_{R/J})
\]
is injective. When $p$ is nilpotent in $R$, this is Katz's
Lemma~1.1.3(2), whose hypotheses are verified for Barsotti--Tate groups in
his proof of Theorem~1.2.1 \cite{Katz1981SerreTate}. Its elementary reason is
that the kernel of reduction in the target formal Lie group is killed by
a fixed power of $p$, whereas the source is $p$-divisible. For completeness,
if $J^2=0$ that kernel is an additive tangent module, hence is killed by
any power of $p$ killing $R$. For general nilpotent $J$, apply this argument
to its successive powers.

For general noetherian $R$, localize at a maximal ideal $\mathfrak m$.
If $p$ is a unit, the finite levels are etale and the assertion is
infinitesimal rigidity for etale schemes. If $p\in\mathfrak m$, reduce
modulo $\mathfrak m^r$ for every $r\geq1$, and use the case just proved.
At each fixed finite level the two coordinate maps differ by entries in
finite locally free modules over $R_{\mathfrak m}$. All these entries lie
in $\bigcap_r\mathfrak m^rR_{\mathfrak m}=0$. Thus every finite-level map
agrees, proving the stated rigidity.

The desired descent assertion is local on $S$, so write $S=\operatorname{Spec}R$.
For each $n$, the functor of group homomorphisms
$G[p^n]\to H[p^n]$ is represented by an affine finitely presented
$S$-scheme: use the affine space of linear maps between their finite
locally free coordinate algebras, and impose the unit, multiplication,
comultiplication, and counit equations. Restriction to the preceding
level gives transition maps. Hence the functor
$X=\underline{\operatorname{Hom}}_S(G,H)$ is represented by their affine
inverse limit. This argument does not assert that $X$ is of finite type.

A homomorphism over $S'$ satisfying the two pullback equalities is a map
$g':S'\to X$. Set
\[
 R_0=\ker\bigl(\Gamma(S',\mathcal O_{S'})
 \rightrightarrows\Gamma(S'',\mathcal O_{S''})\bigr).
\]
The ring $R_0$ is finite as an $R$-module, because it is an $R$-submodule
of the finite module $\Gamma(S',\mathcal O_{S'})$ and $R$ is noetherian.
Also $R\hookrightarrow R_0$. The map $g'$ factors through
$S_0=\operatorname{Spec}R_0$. The map $S'\to S_0$ is proper and
schematically dominant, hence surjective. Two geometric points of $S_0$
over the same geometric point of $S$ lift to points of $S'$; the equalizer
condition identifies them. Consequently $S_0\to S$ is finite, surjective,
and radicial.

Write $X=\operatorname{Spec}C$. Let $B$ be the image of the resulting
$R$-algebra map $C\to R_0$. It is again a finite $R$-module. We have a
finite radicial surjection $Y=\operatorname{Spec}B\to S$, with
$R\hookrightarrow B$, and a closed immersion $Y\hookrightarrow X$.
We claim $\Omega^1_{B/R}=0$. This is a finite $B$-module. The two maps
\[
 B\longrightarrow B\oplus\Omega^1_{B/R},\qquad
 b\longmapsto(b,0),\quad b\longmapsto(b,db)
\]
induce two $p$-divisible-group homomorphisms that agree modulo the
square-zero ideal $\Omega^1_{B/R}$. The test ring is noetherian, so the
rigidity proved above makes them equal. Since $Y\to X$ is a closed
immersion, the two displayed maps agree, and $db=0$ for every $b$.

A finite radicial $R$-algebra with vanishing differentials has, at each
point of $S$, fibre exactly its residue field. The cokernel of $R\to B$
therefore has zero fibre at every point; Nakayama's lemma shows it is zero.
As the map was injective, $B=R$. Thus $g'$ descends to $S$.
Finally, uniqueness follows at every finite level from
$\mathcal O_S\hookrightarrow f_*\mathcal O_{S'}$, since the coordinate
algebras of $G[p^n]$ are finite locally free. This proves the equalizer.
\end{proof}

\begin{proposition}[C-RD2: a limit to infinitesimal rigidity]\label{prop:C-RD2}
\status{Proved; correction to an intermediate claim.}
The functor $\underline{\operatorname{Hom}}(G,H)$ need not be formally
unramified on arbitrary test rings when the base has mixed characteristic.
Consequently ``formally unramified'' in the reading of 100:25 cannot be
used without specifying the allowed thickenings.
\end{proposition}

\begin{proof}
Let $R=\mathbf Z_p\oplus I$ be the square-zero extension with
$I=\mathbf Q_p/\mathbf Z_p$. For $n\geq1$, put
$x_n=p^{-n}+\mathbf Z_p\in I$ and $z_n=1+x_n\in R^\times$.
Then $z_n^{p^n}=1$ and $z_{n+1}^p=z_n$. Sending $1\in\mathbf Z/p^n$
to $z_n$ defines a compatible, nonzero homomorphism
\[
 (\mathbf Q_p/\mathbf Z_p)_R\longrightarrow\mu_{p^\infty,R}
\]
that reduces to zero modulo $I$. The ideal $I$ is square-zero but is not
a finitely generated $\mathbf Z_p$-module. Hence this counterexample is
consistent with Proposition C-RD1: its proof only uses noetherian
thickenings, and passes to a finite schematic image before testing
differentials.
\end{proof}

\subsection{Local systems over geometrically unibranch bases}

\begin{theorem}[C-RD7: effectivity over reduced geometrically unibranch bases]
\label{thm:C-RD7}
\status{Proved.}
Let $S$ be a reduced, locally noetherian, geometrically unibranch scheme.
In particular, $S$ may be any normal locally noetherian scheme.
Every system $A(s)$ of
abelian schemes over $\operatorname{Spec}\mathcal O_{S,s}$, equipped
with the compatible generization isomorphisms in C.4, is induced by an
abelian scheme over $S$, uniquely up to the unique isomorphism
respecting the entire system. No excellence hypothesis is required.
\end{theorem}
\begin{proof}
We first record the exact extension statement used here. If $X$ is a
reduced, geometrically unibranch, locally noetherian integral scheme
with generic point $\eta$,
and $A,B$ are abelian schemes over $X$, then
\[
 \operatorname{Hom}_X(A,B)\longrightarrow
 \operatorname{Hom}_\eta(A_\eta,B_\eta)
\]
is a bijection. Here is a proof which does not require finite
normalization. Work over a noetherian affine open in $X$. The scheme
$H=\operatorname{Hom}_{X\text{-}\mathrm{gp}}(A,B)$ is separated,
unramified, and essentially proper over $X$, meaning that it satisfies
the valuative criterion of properness
\cite[p.64]{Grothendieck1966Hom}. A generic homomorphism spreads over
some dense open $V\subset X$. In the proof of
\cite[Lemma~1.2.1, p.65]{Grothendieck1966Hom}, Grothendieck proves
explicitly that the reduced closure $Z$ of this section in $H$ is
finite and surjective over $X$. This cited closure assertion is
essential: essential properness alone is not being substituted for
properness of the possibly non-quasicompact scheme $H$.

The scheme $Z$ is integral and birational over $X$. On affine opens,
its finite coordinate algebra is a subalgebra of the function field,
and therefore of the integral closure of $\mathcal O_X$. Consequently
the normalization of $X$ maps surjectively to $Z$. Geometric
unibranchness says that this normalization map is a universal
homeomorphism. It follows, by surjectivity after every field extension,
that $Z\to X$ is radicial. As a closed subscheme of $H$, it is also
unramified. A finite radicial unramified morphism is a closed immersion;
its schematic dominance, or simply reducedness and birationality of
$X$, makes that immersion an isomorphism. Thus the section extends.
Separatedness gives uniqueness. Applying the result to a generic
isomorphism and its inverse extends generic isomorphisms as well.

The irreducible components of a locally noetherian geometrically
unibranch scheme are locally disjoint and open. Work on one such
component, which is integral because $S$ is reduced.
Let $\eta$ be its generic point. Since an abelian scheme is of finite
presentation, each $A(s)$ spreads to an abelian scheme $B_s$ over
some open neighbourhood $V_s$ of $s$; the group operations, properness,
smoothness and geometrically connected fibres spread after shrinking.
The given generization maps identify every $(B_s)_\eta$ with the
fixed abelian variety $A(\eta)$. On $V_s\cap V_t$, the resulting
generic isomorphism extends uniquely by the preceding paragraph.
Uniqueness proves the cocycle condition. These isomorphisms glue the
$B_s$ to an abelian scheme $A/S$.

For a point $t$, compare $A$ over $\operatorname{Spec}\mathcal O_{S,t}$
with the prescribed $A(t)$. Both have the stipulated generic
identification with $A(\eta)$; the same Hom extension theorem
therefore gives the unique local isomorphism. Its compatibility with
all generization maps follows from equality at the generic point.
Finally these constructions glue over the locally disjoint components
of $S$.
\end{proof}

\begin{lemma}[Compatible local sections of an unramified scheme]
\label{lem:C-RD8-sections}
\status{Proved under stated hypotheses}
Let $X$ be locally noetherian and let $H\to X$ be separated,
unramified and locally of finite presentation. Suppose a section over
$\operatorname{Spec}\mathcal O_{X,x}$ is given for every $x\in X$,
compatibly with all generizations. Then these sections are induced by
a unique section over $X$.
\end{lemma}
\begin{proof}
Work on a noetherian open of $X$. Spread the section at $x$ to an open
neighbourhood $V_x$. Delete from $V_x$ all irreducible components of
$X$ which do not contain $x$. Every generic point of the remaining
$V_x$ then generalizes $x$, so the spread has the prescribed generic
value. On an overlap, two spreads agree at every generic point.
Separatedness makes them agree on the reduction of the overlap,
because that reduction is schematically dominated by its generic
points. The nilradical is nilpotent locally; formal unramifiedness
makes the two spreads agree on the entire overlap. They glue, and
the same argument, or localization at every point, gives uniqueness.
\end{proof}

\begin{theorem}[C-RD8: reduced bases with a finite geometrically unibranch cover]
\label{thm:C-RD8}
\status{Proved, with the covering hypothesis displayed explicitly.}
Let $S$ be a reduced locally noetherian scheme which, locally on $S$,
admits a finite surjective morphism $S'\to S$ with $S'$ geometrically
unibranch. Then every local system in C.4 is effective, uniquely up to
the unique isomorphism respecting its prescribed identifications.
The normalization of $S$ need not be finite.
\end{theorem}
\begin{proof}
The assertion is local on $S$, with unique compatible gluing, so fix
such a cover. Replace $S'$ by its reduction. Pull the local system
back by assigning to $x\in S'$ the base change of $A(f(x))$ to
$\operatorname{Spec}\mathcal O_{S',x}$. Theorem~\ref{thm:C-RD7}
produces an abelian scheme $B'/S'$ with all these local identifications.

On $D=S'\times_SS'$, these identifications give, at every point $z$
lying over $s\in S$, an isomorphism between the two pullbacks of $B'$:
both are identified with the base change of $A(s)$. These local
isomorphisms are generization-compatible. Their Isom scheme is
separated, unramified and locally of finite presentation, so
Lemma~\ref{lem:C-RD8-sections} gives a global isomorphism on $D$.
The same lemma on the triple fibre product verifies its cocycle
condition. This argument applies even when these fibre products are
nonreduced.

Fix $s\in S$. On $S'\times_S\operatorname{Spec}\mathcal O_{S,s}$,
compare $B'$ with the pullback of $A(s)$. At every local point the
given generization maps and the local identifications of $B'$ give
compatible isomorphisms. The lemma globalizes them on this scheme.
Spread $A(s)$ to an abelian scheme $B_s$ on a neighbourhood $V_s$.
Finite presentation of the schemes and their isomorphism, together
with finiteness of $S'\to S$, spreads the last isomorphism to
\[
 B_s\times_S S'\simeq B'|_{f^{-1}(V_s)}
\]
after shrinking $V_s$. This step spreads one isomorphism of abelian
schemes of finite presentation, and does not spread an infinite BT
system. Over $D\times_S\operatorname{Spec}\mathcal O_{S,s}$ this
isomorphism respects the descent isomorphism just constructed. Since
$D\to S$ is finite, that equality of finitely presented morphisms
also spreads after a further shrinking of $V_s$. Thus the chosen
isomorphism over $f^{-1}(V_s)$ respects the descent isomorphism
throughout $D|_{V_s}$.

It remains to descend the resulting isomorphisms between $B_s$ and
$B_t$ on overlaps. They satisfy the descent equality on $D$, by their
pointwise construction and the preceding lemma. The finite surjection
$S'\to S$ is schematically dominant because $S$ is reduced. Apply
Proposition~\ref{prop:C-RD1} to their BT homomorphisms, obtaining a
unique BT homomorphism over $V_s\cap V_t$. At every generic point
this map is algebraic: faithfully flat descent over the finite
nonzero algebra above that field descends the abelian homomorphism.
Every associated point of the reduced overlap is generic, so
\cite[Proposition~1.2, p.64]{Grothendieck1966Hom} makes the descended
BT map algebraic on the whole overlap. Apply the same argument to its
inverse. Faithfulness verifies both inverse identities and the
cocycle. The $B_s$ therefore glue to an abelian scheme over $S$.

Its local identification with $A(s)$ is already built into $B_s$.
All prescribed generization compatibilities and uniqueness follow
either after the schematically dominant cover or from faithfulness
of the BT functor. This proves the theorem.
\end{proof}

\begin{remark}[The nodal involution does not supply the advertised counterexample]\label{auto:C:63e1d81ef0f5}
\status{Proved under stated hypotheses}

The proposed gluing by $[-1]$ at the two branches of a node has finite
monodromy and is effective: twist the constant abelian scheme by the
corresponding finite etale double cover. This observation refutes that
specific proposed construction, but does not by itself settle C.4 on
all reduced or all nonreduced noetherian bases. In particular, a proof
which shrinks to a normal dense open must justify that such an open
exists; this cannot be silently assumed for an arbitrary nonexcellent
noetherian integral scheme.
\end{remark}

\subsection{Nilpotent invariance of effectivity for a specified local system}

\begin{lemma}[Compatible germs of a coherent module]
\label{lem:C-RD9-coherent}
\status{Proved under stated hypotheses}
Let $X$ be locally noetherian and let $\mathcal F$ be coherent.
A family $m_x\in\mathcal F_x$, compatible with every generization,
is the family of germs of a unique global section of $\mathcal F$.
\end{lemma}
\begin{proof}
Work on a noetherian affine open. Spread $m_x$ to a section on a
neighbourhood $V_x$ and delete the closures of all points in the finite
set $\operatorname{Ass}(\mathcal F)$ which do not generalize $x$.
Every associated point left in $V_x$ then generalizes $x$, so the
section has the prescribed value there. On overlaps, differences
vanish at every associated point of $\mathcal F$. Such a difference
is zero: a nonzero cyclic submodule of $\mathcal F$ would have an
associated prime in $\operatorname{Ass}(\mathcal F)$ at which it
remained nonzero. The local sections therefore glue. Uniqueness is
immediate from their germs.
\end{proof}

\begin{theorem}[C-RD9: nilpotents introduce no new obstruction in C.4]
\label{thm:C-RD9}
\status{Proved; independently audited with the original source.}
Let $S$ be locally noetherian, and fix a local system of abelian
schemes in the sense of C.4. This specified system is effective if
and only if its reduction on $S_{\mathrm{red}}$ is effective.
Consequently every such system is effective whenever
$S_{\mathrm{red}}$ locally admits a finite surjective geometrically
unibranch cover. No assumption that a prime is nilpotent on $S$ is
required.
\end{theorem}
\begin{proof}
Reduction proves necessity. For sufficiency it is enough to handle a
square-zero thickening $S_0\hookrightarrow S$ with coherent ideal
$I$, and a system whose reduction is induced by $A_0/S_0$ with the
specified local identifications. We may first work on an affine open
of $S$.

Consider the Zariski sheaf of isomorphism classes of actual abelian
scheme liftings of $A_0$, with their chosen identification on $S_0$.
It is a sheaf because an automorphism of an abelian scheme reducing
to the identity across a nilpotent ideal is the identity: the Hom
scheme is unramified. Compatible isomorphism classes therefore have
unique descent isomorphisms, and Zariski descent for schemes is
effective. This sheaf is locally nonempty. Indeed every prescribed
local $A(s)$, together with its reduction identification, spreads to
an open neighbourhood by finite presentation.

The square-zero deformation calculation for abelian schemes makes
this sheaf a torsor under the coherent module
\[
 \mathcal F
 =R^1\pi_*\mathcal Hom_{\mathcal O_{A_0}}
       (\Omega^1_{A_0/S_0},\pi^*I)
 =\operatorname{Lie}(A_0)\otimes
   \operatorname{Lie}(A_0^\vee)\otimes I,
 \qquad \pi:A_0\longrightarrow S_0.
\]
The lift-torsor calculation is explicitly recorded in the proof of
\cite[Lemma~24, p.590]{VasiuZink2010Purity}; existence over every affine square-zero thickening is also
proved over an arbitrary ring in
\cite[Chapter~IV, Lemma~2.8.1, pp.148--149]{Messing1972}. Neither
lifting statement imposes a characteristic restriction. Only
the difference-torsor calculation is needed here: local existence is
supplied by the given local system. The displayed tensor identity
uses invariant differentials and
$R^1\pi_*\mathcal O_{A_0}=\operatorname{Lie}(A_0^\vee)$, with
cohomology and base change.

Since $S_0$ is affine, $H^1(S_0,\mathcal F)=0$. The torsor has a
global section, represented by an actual abelian scheme $B/S$.
Compare $B$ at each local point with the prescribed $A(s)$. Their
difference is a well-defined element of $\mathcal F_s$ and these
elements are compatible with all generizations. By
Lemma~\ref{lem:C-RD9-coherent} they are the germs of a section
$m\in\Gamma(S_0,\mathcal F)$. Translate the lifting $B$ by this
section in the torsor of liftings. The resulting abelian scheme has
the required local isomorphisms. They are unique when required to
respect the reduction, and hence obey every prescribed generization
compatibility.

This proves the square-zero assertion on affine opens. On overlaps,
the resulting schemes have the prescribed compatible local
isomorphisms; Lemma~\ref{lem:C-RD8-sections}, applied to their Isom
scheme, globalizes those isomorphisms and verifies their cocycle.
Thus the assertion holds on all locally noetherian $S$.

Locally the nilradical is nilpotent. Apply the square-zero result
successively to its finite filtration by powers, and then glue the
constructions by the same uniqueness. This proves sufficiency for
$S_{\mathrm{red}}\hookrightarrow S$. The last assertion follows
from Theorem~\ref{thm:C-RD8}. Every lifting and gluing in this proof
is at a finite nilpotent stage and is an actual scheme; no assertion
that an arbitrary infinite formal deformation algebraizes is used.
\end{proof}

\begin{remark}[A genuinely nonreduced failure]\label{auto:C:24b235bcb6c7}
\status{Proved under stated hypotheses}

For the ring $R$ and local system of Theorem~\ref{thm:C-RD10}, base
change each local scheme and its generic identification along
$R\to R[\epsilon]/(\epsilon^2)$. This is a local system on a
nonreduced noetherian base. Any global realization would reduce to
the impossible realization over $R$, so it also fails to be effective.
\end{remark}

\subsection{An explicit reduced counterexample to unrestricted C.4}

\begin{theorem}[C-RD10: infinitely many nodes obstruct the local-system lemma]
\label{thm:C-RD10}
\status{Proved; independently audited.}
There is a countable, one-dimensional, noetherian integral
$\mathbf Q$-algebra $R$, all of whose closed local rings are ordinary
split nodes, and a generization-compatible family of elliptic schemes
on $\operatorname{Spec}\mathcal O_{\operatorname{Spec}R,s}$ which does
not come from an abelian scheme over $\operatorname{Spec}R$.
Consequently C.4 is false for general reduced bases and hence for
arbitrary locally noetherian bases. The counterexample is explicit
and uses no cardinality diagonalization.
\end{theorem}
\begin{proof}
\emph{The ring.} Let $t_1,t_2,\ldots$ be algebraically independent over
$\mathbf Q$, set $x_i=t_i^2-1$ and $y_i=t_i(t_i^2-1)$, and put
\[
 C_i=\mathbf Q[x_i,y_i]
   \simeq\mathbf Q[X,Y]/(Y^2-X^2(X+1)),\qquad
 R_0=\bigotimes_{i\geq1,\mathbf Q} C_i\subset\mathbf Q[t_1,t_2,\ldots].
\]
Write $P_i=(x_i,y_i)R_0$, set
\[
 \Sigma=R_0\setminus\bigcup_{i\geq1}P_i,
 \qquad R=\Sigma^{-1}R_0,
 \qquad \mathfrak m_i=P_iR.
\]
Each $P_i$ is prime, since its quotient is the tensor product of the
other geometrically integral nodal rings. Thus $\Sigma$ is
multiplicative. A nonzero element of $R_0$ involves only finitely many
indices, and belongs to none of the $P_i$ outside those indices.

We verify all the finiteness properties used below. If a nonzero prime
$Q\subset R_0$ avoids $\Sigma$, choose $0\ne f\in Q$. The set
$I=\{i:f\in P_i\}$ is finite. If $Q$ were contained in none of these
$P_i$, finite prime avoidance would give
$g\in Q\setminus\bigcup_{i\in I}P_i$. Let $J$ contain every index
occurring in $f$ and $g$. For $i\in I$, $f+cg\notin P_i$ whenever
$0\ne c\in\mathbf Q$. For each $i\in J\setminus I$ there is at most
one scalar $c$ with $f+cg\in P_i$. Choose a nonzero rational number
avoiding these finitely many exceptions and also ensuring $f+cg\ne0$.
For $i\notin J$, independence of $t_i$ gives $f+cg\notin P_i$.
Then $f+cg\in Q\cap\Sigma$, a contradiction. Every prime of $R$
is therefore contained in some $\mathfrak m_i$. These ideals are
pairwise incomparable: $x_i\in P_i\setminus P_j$ for $i\ne j$, a
property unchanged by localization. Hence they are exactly the
maximal ideals of $R$.

Let $L_i=\mathbf Q(t_j:j\ne i)$. Localizing first in the nonzero
coefficients from the other factors gives
\[
 R_{\mathfrak m_i}
   \simeq L_i[X,Y]_{(X,Y)}/(Y^2-X^2(X+1)).                 \tag{C-RD10.1}
\]
These are noetherian one-dimensional local domains. Every nonzero
$a\in R$ belongs to only finitely many maximal ideals. If $I\subset R$
is a nonzero ideal, choose $0\ne a\in I$. At each of those finitely
many maximal ideals containing $a$, choose finitely many elements of
$I$ generating its localization. Together with $a$ they generate an
ideal $J\subset I$ whose localization equals $I$ at every maximal
ideal; therefore $I/J=0$. Thus every ideal is finitely generated,
proving that $R$ is noetherian. Formula (C-RD10.1) proves $\dim R=1$.
Its points are the generic point $\eta$ and the $\mathfrak m_i$.
The tangent cone $Y^2-X^2$ has two distinct rational branches.

The element $t_i=y_i/x_i$ lies in $F=\operatorname{Frac}R$ and satisfies
$t_i^2=1+x_i$, but does not belong to $R_{\mathfrak m_i}$. Indeed,
every rational function in (C-RD10.1) regular at the node has equal
values at the two points $t_i=1$ and $t_i=-1$ of its normalization:
write it as $a(X,Y)/b(X,Y)$ with $b(0,0)\ne0$ and evaluate. The
function $t_i$ has the distinct values $1$ and $-1$.
This also proves that the normalization of $R$ is not finite. If it
were finite, some nonzero $d\in R$ would multiply it into $R$; choose
$i$ with $d\notin\mathfrak m_i$. The integral element $t_i$ would
then belong to $R_{\mathfrak m_i}$, a contradiction.

\emph{The local elliptic schemes.} Fix
\[
 E/\mathbf Q:\quad Y^2=X^3+X+1,
 \qquad \Delta(E)=-496.
\]
At the generic point use $E_F$. At $\mathfrak m_i$, put $d_i=1+x_i$
and take
\[
 A(\mathfrak m_i):\quad
 Y^2=X^3+d_i^2X+d_i^3
 \quad\text{over }R_{\mathfrak m_i}.                    \tag{C-RD10.2}
\]
The element $d_i$ is a unit, and the discriminant is $-496d_i^6$,
so this is an elliptic scheme. Its generic fibre is identified with
$E_F$ by $(X,Y)\mapsto(X/t_i^2,Y/t_i^3)$. Since the only nonidentity
generizations on $\operatorname{Spec}R$ run from a closed point to
$\eta$, these identifications define the required local system.

However (C-RD10.2) is not isomorphic to the constant $E$ over
$R_{\mathfrak m_i}$. Since $2$ and $3$ are invertible, an isomorphism
of short Weierstrass elliptic schemes preserving the zero section is
of the form $X\mapsto u^2X$, $Y\mapsto u^3Y$, with $u$ a unit;
comparison of the absent $X^2$, $XY$ and $Y$ terms eliminates all
translations. Comparing the two nonzero coefficients gives
$u^4=d_i^2$ and $u^6=d_i^3$, hence $u^2=d_i$. In the field $F$ this
forces $u=\pm t_i$, which are not elements of $R_{\mathfrak m_i}$.

\emph{Failure of effectivity.} If an abelian scheme $A/R$ realized the
system, its generic fibre would be identified with $E_F$. Finite
presentation spreads this generic isomorphism and its inverse to
some principal dense open $D(f)$, with $0\ne f\in R$. But $f$ lies
in only finitely many $\mathfrak m_i$. For every other $i$, one would
have
\[
 A(\mathfrak m_i)\simeq A\times_RR_{\mathfrak m_i}
   \simeq E\times_{\mathbf Q}R_{\mathfrak m_i},
\]
contradicting the preceding paragraph. This proves noneffectivity. The exact polynomial identities and
Weierstrass-change coefficients are also verified in
\texttt{verify\_nodal\_twist.py}, with its saved JSON certificate.
\end{proof}

\begin{remark}[Scope of the example]\label{auto:C:3a7ebde95252}
\status{Proved under stated hypotheses}

The infinite tensor-product localization follows the construction
attributed to \cite[Proposition~1 and Example~1]{Hochster1973Loci},
with a nodal factor in place of the cuspidal one. The original
Hochster article has not yet been obtained: author bibliographical
metadata and the Stacks reference were checked, and the proposition
was located in a secondary quotation. The preceding proof independently
checks noetherianity and every property used in the counterexample.
It does not contradict Theorem~\ref{thm:C-RD8}: its conclusion in fact
shows that $\operatorname{Spec}R$ has no open cover by schemes admitting
finite surjective geometrically unibranch covers. The single-node
$[-1]$ proposal in Lead C remains effective; the obstruction here is
that infinitely many nontrivial local twists cannot be supported on
the complement of any dense open.
\end{remark}


\subsection{The exact additional condition across a nilpotent thickening}

\begin{lemma}[Serre--Tate across a finite nilpotent thickening of a complete base]
\label{lem:C-RD12-complete}
\status{Proved under stated hypotheses}
Let $R$ be a noetherian $p$-adically complete ring with $p$ in its
Jacobson radical, let $J\subset R$ be nilpotent, and put $R_0=R/J$.
Fix an abelian scheme $A_0/R_0$, a BT group $M/R$, and an
isomorphism $A_0[p^\infty]\simeq M_{R_0}$. There is an abelian scheme
$A/R$ with these prescribed reductions, unique up to a unique
isomorphism respecting the prescribed reduction and BT comparison.
No polarization is required. This assertion concerns a finite
nilpotent thickening of an already algebraic abelian scheme, not
algebraization of an arbitrary infinite formal abelian deformation.
\end{lemma}
\begin{proof}
First suppose $J^2=0$. Actual abelian liftings of $A_0$ are locally
nonempty and form a torsor under
\[
 \mathcal F=\operatorname{Lie}(A_0)\otimes
             \operatorname{Lie}(A_0^\vee)\otimes J.
\]
Here and below a lifting includes the identification of its reduction.
The actual-lifting torsor, and its triviality over a local noetherian
base, are recorded explicitly in the proof of
\cite[Lemma~24, p.590]{VasiuZink2010Purity}; existence over every affine square-zero thickening is also
proved in \cite[Chapter~IV, Lemma~2.8.1, pp.148--149]{Messing1972}. Spread local liftings,
and use $H^1(\operatorname{Spec}R_0,\mathcal F)=0$ to choose an actual
lifting $B/R$. There are no infinitesimal automorphisms of an abelian
scheme as a group scheme, so this torsor of isomorphism classes is
an honest Zariski sheaf and its gluing is effective.

For each $r\geq1$, apply Serre--Tate to
\[
 R/(J+p^rR)\longleftarrow R/p^rR.
\]
It gives the unique lifting $A_r$ having BT group $M\bmod p^r$ and
the specified reduction of $A_0$. Relative to $B\bmod p^r$, its
difference is an element
\[
 d_r\in\Gamma\!\left(\operatorname{Spec}R_0,
   \operatorname{Lie}(A_0)\otimes\operatorname{Lie}(A_0^\vee)
        \otimes J/(J\cap p^rR)\right).
\]
These elements are compatible. The induced filtration
$J\cap p^rR$ is cofinal with the $p$-adic filtration on $J$ by
Artin--Rees. Since $J$ and the displayed coherent module are complete,
$(d_r)_r$ comes from a unique $d\in\Gamma(\mathcal F)$.
Translate the actual lifting $B$ by $d$. The resulting actual abelian
scheme $A/R$ has all the prescribed reductions $A_r$.

For each fixed $m$, the compatible isomorphisms
$A[p^m]\bmod p^r\simeq M[p^m]\bmod p^r$ lift uniquely to an
isomorphism over $R$: their coordinate algebras are finite projective
$R$-modules and hence complete. The Hopf identities, inverses and
compatibilities in $m$ are finite equalities and pass to the limit.
This gives $A[p^\infty]\simeq M$. If two liftings have the same data,
their difference in $\Gamma(\mathcal F)$ is zero modulo every
$J\cap p^rR$, hence is zero. Their reduction-preserving isomorphism
is unique by unramifiedness of the abelian Hom functor.

Finally induct along the finite filtration by powers of $J$.
Every step is a square-zero lifting of an actual abelian scheme,
which proves existence and uniqueness as stated.
\end{proof}

\begin{theorem}[C-RD12: exact nilpotent correction to the extension assertion]
\label{thm:C-RD12}
\status{Proved; independently audited against the original sources.}
Let $S$ be locally noetherian, let $T\subset V(p)$ be closed, and put
$U=S\setminus T$. Let $S_0\hookrightarrow S$ be a nilpotent closed
immersion. Fix a triple $(A_U,M,\varphi)$ as in C.1 and a specified
realization $A_0/S_0$ of its reduction.

Then this triple has a realization $A/S$ inducing the specified
$A_0$ if and only if the following condition holds on an affine open
cover of $S$. For $V=\operatorname{Spec}R$, choose a finitely generated
ideal $I$ containing $p$ with $T\cap V=V(I)$, and set
\[
 \widehat V=\operatorname{Spec}\widehat R_I,
 \qquad \widehat U=\widehat V\times_S U.
\]
Let $\widehat A/\widehat V$ be the unique algebraic lifting of
$A_0\times_{S_0}(\widehat V\times_SS_0)$ with BT group
$M_{\widehat V}$ supplied by Lemma~\ref{lem:C-RD12-complete}.
There must be an isomorphism
\[
 \widehat A|_{\widehat U}\ \simeq\ A_U|_{\widehat U}
                                                        \tag{C-RD12.1}
\]
which induces the specified reduction and the specified BT
identification $\varphi$. Thus, once the reduced triple is realized,
the punctured-completion comparison is an exact condition, with no
additional polarization, excellence, or formal-algebraization
hypothesis. It is an additional condition: it fails for
Proposition~\ref{prop:C-RD5}.
\end{theorem}
\begin{proof}
The completion $\widehat R_I$ is noetherian and $p$-adically complete.
Indeed a $p$-adically Cauchy sequence is $I$-adically Cauchy, and the
ideals $p^r\widehat R_I$ are $I$-adically closed; separatedness follows
from the Krull intersection theorem, since $p\in I$ lies in the
Jacobson radical. The kernel of
$\widehat R_I\to\widehat R_I/J\widehat R_I$, where $J$ defines
$S_0$, is nilpotent. Lemma~\ref{lem:C-RD12-complete} therefore applies
to this kernel; it is not necessary that it equal the entire
nilradical of $\widehat R_I$.

Necessity is base change and the uniqueness in that lemma. For
sufficiency, first suppose $J^2=0$ and work on $V=\operatorname{Spec}R$.
As in the lemma, choose an actual lifting $B/V$ of $A_0|_{V\cap S_0}$.
Let $\mathcal F$ be its coherent difference module. The schemes
$A_U|_{U\cap V}$ and $\widehat A$ have respective differences
\[
 d_U\in\Gamma(U\cap V,\mathcal F),\qquad
 d_{\widehat V}\in\Gamma(\widehat V,\mathcal F_{\widehat V})
\]
from the corresponding restrictions of $B$. Condition (C-RD12.1)
says exactly that these sections agree on $\widehat U$.
The quasi-coherent gluing equivalence of Ferrand--Raynaud, in the
precise form \cite[Theorem~0.3, p.561]{MoretBailly1996Descent}, gives
\[
 \Gamma(V,\mathcal F)\ \simeq\
 \Gamma(U\cap V,\mathcal F)
 \mathop{\times}_{\Gamma(\widehat U,\mathcal F)}
 \Gamma(\widehat V,\mathcal F_{\widehat V}).
\]
Translate $B$ by the resulting section. This produces an actual
abelian scheme $A/V$ with the required isomorphisms over $U\cap V$
and $\widehat V$.

Its BT group has the required comparison with $M$. At each finite
level, descend the compatible Hopf-algebra isomorphisms over
$U\cap V$ and $\widehat V$ by the same quasi-coherent equivalence.
Their group laws and compatibility between levels descend as
identities. Thus $A[p^\infty]\simeq M$, with the prescribed
identification over $U\cap V$.

The construction is unique among liftings with the specified
reduction and these restrictions: its difference from another is a
section of $\mathcal F$ vanishing on $U\cap V$ and $\widehat V$,
hence zero by the displayed equalizer. On overlaps the same argument,
after restriction to smaller affine opens, provides the unique
reduction-preserving isomorphisms, so the schemes glue. Induction
along the finite powers of $J$ gives the general nilpotent case.
At each stage the two prescribed restrictions are the reductions of
$A_U$ and $\widehat A$, so (C-RD12.1) provides the required compatible
differences. This proves sufficiency globally.
\end{proof}

\begin{remark}[What this criterion does and does not remove]\label{auto:C:35345d1cc7ea}
\status{Proved under stated hypotheses}

Taking $S_0=S_{\mathrm{red}}$ separates the essential-image problem
into the reduced realization problem and an exact additional
nilpotent comparison. Theorem~\ref{thm:C-RD11} shows that the first
problem can already fail under a) and b). Consequently one cannot
replace it by the assumption ``$S$ is reduced''. The theorem above
settles the extra condition for nonreduced bases once a reduced
realization is specified; it does not assert automatic effectivity
of a general unpolarized formal abelian scheme.
\end{remark}


\subsection{An absolute essential-image criterion}

The criterion below is an existence statement about explicit data on
the closed locus, its infinitesimal neighbourhoods, the completion,
and the punctured completion. In particular it does not quantify over
an abelian scheme on the original base. Its two effectivity conditions
are retained because unpolarized formal algebraization and descent of
schemes cannot be taken for granted.

\begin{theorem}[C-RD13: the exact image by formal models and compatible charts]
\label{thm:C-RD13}
\status{Proved; independent root audit passed.}
Let $S$ be a locally noetherian scheme, let $T\subset V(p)$ be a closed subset,
put $U=S\setminus T$, and fix a triple $(A_U,M,\varphi)$ as in C.1.
Choose a coherent ideal $\mathcal I\subset\mathcal O_S$ containing
$p$, with $V(\mathcal I)=T$, and put
\[
 T_r=V(\mathcal I^r)\quad(r\geq1).
\]
The triple is in the essential image of $F$ if and only if there
exist the following data.

\begin{enumerate}
\item An abelian scheme $A_1/T_1$ and a specified isomorphism
$\eta_1:A_1[p^\infty]\simeq M_{T_1}$.
For $r\geq1$, let $(A_r,\eta_r)$ be the compatible Serre--Tate
liftings of this datum to $T_r$, using $M_{T_r}$. These are determined
up to their unique reduction-preserving compatible isomorphisms.

\item An affine open cover $S=\bigcup_i V_i$, where
$V_i=\Spec R_i$, and, for each $i$, an actual abelian scheme
$B_i/\widehat V_i$, with
\[
 \widehat V_i=\Spec\widehat{R_i}_{I_i},\qquad
 I_i=\mathcal I(V_i),\qquad
 \widehat U_i=\widehat V_i\times_SU,
\]
and compatible identifications
$B_i\times_{\widehat V_i}(T_r\cap V_i)\simeq A_r|_{T_r\cap V_i}$.
Thus $B_i$ is an algebraization of a specified formal object, rather
than an unspecified good model. Its formal BT identifications
uniquely give an actual isomorphism
\[
 \widehat\eta_i:B_i[p^\infty]\simeq M_{\widehat V_i}.
                                                        \tag{C-RD13.1}
\]

\item A group-scheme isomorphism
\[
 \beta_i:B_i|_{\widehat U_i}\ \simeq\ A_U|_{\widehat U_i}
                                                        \tag{C-RD13.2}
\]
whose map on BT groups is
$\varphi|_{\widehat U_i}\circ\widehat\eta_i|_{\widehat U_i}$.
The condition requires an actual abelian-scheme isomorphism on this
overlap, and not merely an isomorphism of its BT groups.

\item For each $i$, finite open covers
\[
 B_i=\bigcup_a W_{ia},\qquad
 A_U|_{U\cap V_i}=\bigcup_a Z_{ia},
                                                        \tag{C-RD13.3}
\]
where $W_{ia}\to\widehat V_i$ and $Z_{ia}\to U\cap V_i$ are
quasi-affine, with a common finite index set, and $\beta_i$ identifies
$W_{ia}|_{\widehat U_i}$ with $Z_{ia}|_{\widehat U_i}$ for every $a$.
Empty members are allowed.
\end{enumerate}
The ideal and affine cover may be chosen arbitrarily: existence for
one such choice is equivalent to existence for every choice. No
excellence assumption or finite-normalization assumption occurs.
The statement therefore applies in particular under a) and b) of C.1,
and on reduced and nonreduced bases alike.
\end{theorem}
\begin{proof}
We first justify the objects and comparisons in the statement.
Since $p\in\mathcal I$, the element $p$ is nilpotent on every $T_r$.
Each immersion $T_r\hookrightarrow T_{r+1}$ is nilpotent. The
Serre--Tate equivalence
\cite[Theorem~1.2.1]{Katz1981SerreTate} consequently constructs the
compatible liftings $A_r$, including their BT identifications, from
$(A_1,\eta_1)$. This construction works globally on the closed
schemes by the same full faithfulness and gluing. Given $B_i$ in
condition 2, for every fixed $n$ the coordinate algebras of
$B_i[p^n]$ and $M[p^n]_{\widehat V_i}$ are finite projective modules
over the complete noetherian ring $\widehat R_i$. The compatible
isomorphisms modulo $I_i^r$ therefore give a unique Hopf-algebra
isomorphism over $\widehat R_i$. Passing through all $n$ proves
(C-RD13.1), including compatibility between levels.

Necessity is direct but also checks the atlas requirement. If $A/S$
realizes the triple, take $A_1=A|_{T_1}$ and
$B_i=A\times_S\widehat V_i$. Serre--Tate identifies their formal
reductions with the stated $A_r$, and the given identification with
$A_U$ supplies $\beta_i$. Choose a finite affine open cover of
$A|_{V_i}$, which is possible because it is a noetherian
quasi-compact scheme. An affine scheme mapping to the affine $V_i$
is affine over $V_i$. Its two base changes are therefore affine,
hence quasi-affine, over $\widehat V_i$ and $U\cap V_i$. These are
the paired opens in (C-RD13.3). This proves necessity for every
choice of $\mathcal I$ and of the affine cover.

For sufficiency first work on one affine $V=\Spec R$, omitting $i$.
The completion morphism $\widehat V\to V$ is flat, its inverse image
of $T_1\cap V$ is that same scheme, and
\[
             (U\cap V)\amalg\widehat V\longrightarrow V
                                                        \tag{C-RD13.4}
\]
is faithfully flat and quasi-compact. Apply the quasi-affine gluing
equivalence of \cite[Theorem~1.1, p.565]{MoretBailly1996Descent}
to each paired open $(W_a,Z_a)$ with its restricted $\beta$.
It gives a quasi-affine $V$-scheme $X_a$ with those two prescribed
restrictions. Intersections $W_a\cap W_b$ and $Z_a\cap Z_b$ are
again relatively quasi-affine and are paired by $\beta$. The same
theorem gives $X_{ab}$ and morphisms
\[
                         X_a\longleftarrow X_{ab}\longrightarrow X_b.
\]
These morphisms are open immersions: that property can be checked
after the faithfully flat cover (C-RD13.4), where they are the
original open inclusions. By full faithfulness, their transition
isomorphisms satisfy the cocycle condition, since the condition
holds on each of the two restrictions. Gluing the $X_a$ along these
open subschemes gives a $V$-scheme $A_V$ with specified isomorphisms
\[
 A_V|_{U\cap V}\simeq A_U|_{U\cap V},\qquad
 A_V\times_V\widehat V\simeq B.
                                                        \tag{C-RD13.5}
\]
Their comparison is exactly $\beta$. Finite presentation, smoothness,
properness and geometrically connected fibres are faithfully flat
local properties on the base; they hold on (C-RD13.4), and hence
hold for $A_V/V$.

The group structure also descends without an extra effectiveness
assumption. The two additions give compatible morphisms on the two
restrictions of $A_V\times_VA_V$. The full faithfulness for schemes
in \cite[Theorem~1.2, p.565]{MoretBailly1996Descent} gives a unique
addition on $A_V$. Apply the same assertion to the identity and
inverse. The group identities hold after (C-RD13.4), hence hold on
$A_V$. Thus $A_V$ is an abelian scheme. At every finite level, the
two isomorphisms $A_V[p^n]\simeq M[p^n]$ on this cover agree on
$\widehat U$ by condition 3. The quasi-coherent gluing equivalence
\cite[Theorem~0.3, p.561]{MoretBailly1996Descent}, applied to the
finite coordinate algebras, glues them to an isomorphism over $V$.
All Hopf identities and all transitions in $n$ glue as equalities.
We have therefore realized the complete triple over $V$.

We spell out why these affine constructions glue, even though their
completed models were chosen separately. Take an affine open
$V=\Spec R$ inside $V_i\cap V_j$ and complete it along $\mathcal I$.
The base changes of $B_i,B_j$ to $\widehat V$ are proper schemes
whose formal completions are both the prescribed tower $A_r$ on
$T_r\cap V$. Full faithfulness of formal completion for proper
schemes over a complete noetherian ring gives the unique
isomorphism between these two models inducing those formal
identifications; see \cite[Theorem~5.4.1, pp.156--157]{EGAIII1}.
It preserves the group structures and their BT comparisons, since
both assertions are equalities that hold on every formal reduction.
On $\widehat V\times_SU$, it agrees with the comparison through
$A_U$: the two isomorphisms induce the same BT isomorphism, and the
map from homomorphisms of abelian schemes to homomorphisms of their
BT groups is faithful. The latter faithfulness follows on geometric
fibres from the schematic density of all $p$-power torsion, and
then on the noetherian base from unramifiedness of the abelian Hom
functor. Moret--Bailly's Theorem~1.2 now gives the unique
isomorphism between $A_{V_i}|_V$ and $A_{V_j}|_V$ with the given
open and completed restrictions. These isomorphisms agree on
smaller affine opens and satisfy the cocycle condition by the same
full faithfulness. They therefore glue the $A_{V_i}$ to an abelian
scheme on $S$, and their BT identifications glue as well.

This proves sufficiency. Combined with necessity, it proves the
independence of the chosen ideal and cover asserted in the theorem.
\end{proof}

\begin{corollary}[C-RD14: the exact elliptic criterion]
\label{cor:C-RD14}
\status{Proved; independent root audit passed.}
Suppose $A_U$ has relative dimension one and $M$ has height two. In Theorem~\ref{thm:C-RD13}, conditions 2 and 4 may be
removed. More explicitly, choose an elliptic scheme $A_1/T_1$ with
$A_1[p^\infty]\simeq M_{T_1}$. On each affine completion the
Serre--Tate tower has a unique algebraization $B_i$, and the triple
extends if and only if these $B_i$ admit the actual overlap
isomorphisms (C-RD13.2) respecting the specified BT comparisons.
There is no additional ample-class or scheme-atlas condition in
relative dimension one.
\end{corollary}
\begin{proof}
The zero sections of the elliptic schemes $A_r$ define compatible
relatively ample invertible sheaves
$\mathcal L_r=\mathcal O_{A_r}(3e_r)$. On the affine closed base
$T_1\cap V_i$ their restrictions are ample as invertible sheaves.
Grothendieck's algebraization theorem
\cite[Theorem~5.4.5, pp.157--158]{EGAIII1} gives a projective scheme $B_i$ over
$\widehat V_i$ with this formal completion and this ample line
bundle. The group operations algebraize by proper formal full
faithfulness. At a point over $V(I_i)$, all the quotients
$B_i\bmod I_i^r$ are flat over $\widehat R_i/I_i^r$. The local
formal criterion for flatness, applied to these noetherian local
rings, makes $B_i$ flat at that point. Its fibre is the smooth
fibre of the elliptic scheme $A_1$, so the finite-presentation
smoothness criterion makes $B_i$ smooth there. The complement
of the smooth locus has closed image by properness, and every
nonempty closed subset of the complete affine
base meets $V(I_i)$ because $I_i$ is in the Jacobson radical.
Geometric connectedness then follows from the finite etale
component scheme of a proper smooth group. Thus $B_i$ is elliptic.
The finite-level completeness argument supplies its prescribed BT
comparison, and formal full faithfulness gives its uniqueness.

Any group isomorphism $\beta_i$ preserves the zero sections and
therefore identifies the canonical line bundles
$\mathcal O_{B_i}(3e)$ and $\mathcal O_{A_U}(3e)$ on the overlap.
Apply \cite[Theorem~1.1, the QPROJ case, p.565]{MoretBailly1996Descent}
to these paired schemes with their relatively ample line bundles.
This directly produces a scheme on $V_i$, so no choice of a
compatible quasi-affine cover is necessary. The properness, group
structure, BT identification and global gluing follow as in the
proof of Theorem~\ref{thm:C-RD13}. Necessity remains restriction.
\end{proof}

\begin{remark}[Relation to the nilpotent correction]\label{auto:C:cfff6b2c31f1}
\status{Proved under stated hypotheses}

The general criterion retains two distinct effectivity checks:
algebraization of the formal Serre--Tate object and a compatible
scheme atlas on the open-completion pair. They are explicit
conditions on the given triple and auxiliary data away from $S$;
neither says that the triple already extends over $S$.
Theorem~\ref{thm:C-RD12} is stronger when a reduced realization is
already specified: it proves that both checks can then be removed,
leaving exactly the additional completed-overlap isomorphism.
Corollary~\ref{cor:C-RD14} gives a second situation where both checks
disappear. The reduced counterexample C-RD11 and the square-zero
counterexample C-RD5 show that the remaining comparison cannot be
replaced by reducedness of the local rings.
\end{remark}

